% Results. Brief: briefs/results.md (subsection order set at PH-W8.3: regime map, dissolution, carbonate
% control and the oxidation map, phosphates, nucleation, required precision, baselines). Claims carried:
% the entries of ../claims.yaml with section results. Prose: PH-W8.3, in the regime frame. Each subsection
% opens with the physical question and its physical reading, then the numbers; every number is a macro
% from a committed output file. One sentence per line; each ends with a comment naming its ledger id. Sentences of
% at most 35 words (PH-W8.13).
% carbonate control (frozen): the control row and its paragraph are taken unchanged from the carbonate screen.
\section{Results}\label{sec:results}

\subsection{The regime map}\label{sec:res-regime}
% (1) The physical question and the map that answers it.
A mineral on the collector is asked two separate questions. % CL-S13
Does it stay undissolved in the alkaline water it meets during processing, and does it stay unoxidised at the potential the cathode sets in the cell? % CL-S13
The regime map answers both per phase, per stage and per cutoff class (Table~\ref{tab:regime}); the per-pH tables of the acid-adjusted state are in \nacreSIref{si:regime}{Regime table detail and later reruns}. % CL-S13, CL-D04
A phase that is not thermodynamically excluded (open) has no lower bound on its onset, and neither word implies feasibility or a pass. % CL-D21
Every tier-one exclusion in this paper holds for a \ce{LiClO4}-saturated sink (solid-salt reference state), marked \nacreSinkMark\ from here on. % CL-M17
The sink is a perchlorate electrolyte, because \ce{LiClO4} is the fluorine-free lithium salt with tabulated thermochemistry; nothing is stated for bis(oxalato)borate or any other anion. % CL-M17, CL-L05
These exclusions are thermodynamic ceilings: kinetic stabilisation is not excluded, and the potential hold of Section~\ref{sec:predictions} measures it. % CL-L09, CL-Q01
Calcite meets the dissolution condition over a named region of the alkaline grid, and carbonate is thermodynamically excluded\nacreExcludedMark\nacreSinkMark\ above \NumCarbonateCrossover. % CL-S13
Bone mineral shares that ceiling; at the higher cutoff classes only the condensed calcium phosphates remain not thermodynamically excluded, and the aluminium phosphates are undetermined. % CL-S13
No phase passes the oxidation condition in this version. % CL-S16
$E_\mathrm{lb}$ bounds the onset only if $\Delta G_\mathrm{tr}^\mathrm{rel}\ge0$, which no measured datum tests, so the map can exclude a phase but not admit one. % CL-S16

% Table 6 of the re-audit (break-even ratios) is set before the regime-table floats, so it leaves the float
% queue near Section 4.3, which cites it; behind the three full-width regime parts it floated four pages on.
\begin{table}[t]
\caption{Break-even ratio $r^*$ of the saturated-sink exclusions\nacreSinkMark, per phase and cutoff: the headroom $E_\mathrm{c}-E_\mathrm{ub}-m$ of the upper bound and its margin below the cutoff, and the mean-activity ratio of saturated to actual \ce{LiClO4} at which the sink term $m_\mathrm{sink}$ uses up that headroom. A negative headroom is not an exclusion and has no ratio.}
\label{tab:breakeven}
\nacreBreakevenTable
\end{table}

% The regime table is taller than a page: it is set as consecutive full-width floats, one per part, the
% later parts continuing the table number (cp6 F1; parts and notes from the generator).
% Generated by `nacre preprint regime-table` from the hash-pinned ladder and the production outputs. Do not edit.
% The file defines four macros: \nacreRegimeTable (Results), \nacreRegimePerPh, \nacreRegimeReruns and \nacreAlSaltTable (Supplementary Information).
% The Al-salt table is generated from figures/out/tables/al_salt_ladder.csv (sha256 44635488428bc08439030849796d2eaee34ddf0c6a7d01b5e37741d5c41d44e5).
% Sourced dose key for the phosphoric-acid slurry: p_total_slurry_h3po4_window (a gap). The ratio-grid run is session PH-W4.6.
\newcommand{\nacreRegimeTable}{%
\begingroup\footnotesize\setlength{\tabcolsep}{3pt}
\begin{tabular}{>{\raggedright\arraybackslash}p{0.18\textwidth}l>{\raggedright\arraybackslash}p{0.17\textwidth}>{\raggedright\arraybackslash}p{0.25\textwidth}>{\raggedright\arraybackslash}p{0.11\textwidth}>{\raggedright\arraybackslash}p{0.11\textwidth}>{\raggedright\arraybackslash}p{0.11\textwidth}}
\toprule
row & stage & processing stage (alkaline slurry), $K_2$ & acid-adjusted comparator, $K_2$: sourced slurries pH~\NumAcidSourcedLo{} to \NumAcidSourcedHi{}; grid pH~\NumRegAcidPhLow{} to \NumRegAcidPhHigh{} (spans the Al passivation window) & \NumClassLo{} nominal, $K_1$; the \NumClassSensitivity{} sensitivity line gives the same entry in every row & \NumClassMidLow{} to \NumClassMidHigh{}, $K_1$ & \NumClassHiLow{} to \NumClassHiHigh{}, $K_1$ \\
\midrule
\multicolumn{7}{l}{\emph{carbonate: the polymorph rows are $K_2$ (dissolution), the control row is $K_1$ (oxidation)}} \\
\addlinespace[2pt]
calcite (suitable against dissolution, a necessary condition, in [R-slurry $K_2$ grid \NumRegGridId{}: $\mathrm{Ca_T}\ge$ \NumRegWindowFullLoCalcite{}, \NumRegSlurryPassFullCalcite{} cells, plus \NumRegSlurryPassExtraCalcite{} cells at \NumRegWindowLoCalcite{}]) & slurry & suitable against dissolution (\NumRegSlurryPassCalcite{} of \NumRegSlurryCellsCalcite{} cells: $\mathrm{Ca_T}\ge$ \NumRegWindowFullLoCalcite{}, \NumRegSlurryPassFullCalcite{} cells, plus \NumRegSlurryPassExtraCalcite{} cells at \NumRegWindowLoCalcite{}) & excluded against dissolution\nacreDissolutionMark{} below pH~\NumRegAcidFirstCalcite{} (lower passivation window, below the sourced range); suitable against dissolution from pH~\NumRegAcidFirstCalcite{} up (\NumRegAcidPassCalcite{} of \NumRegAcidCellsAboveCalcite{} cells from that pH up; lowest passing $\mathrm{Ca_T}$ \NumRegAcidMinCaCalcite{} at pH~\NumRegAcidFirstCalcite{} on the crossed grid; at or above \NumRegHighDose{}) & -- & -- & -- \\
aragonite (dissolution stability only; formation unsupported at slurry pH and room temperature) & slurry & suitable against dissolution (\NumRegSlurryPassAragonite{} of \NumRegSlurryCellsAragonite{} cells; $\mathrm{Ca_T}$ \NumRegWindowLoAragonite{} to \NumRegWindowHiAragonite{}) & excluded against dissolution\nacreDissolutionMark{} below pH~\NumRegAcidFirstAragonite{} (lower passivation window, below the sourced range); suitable against dissolution from pH~\NumRegAcidFirstAragonite{} up (\NumRegAcidPassAragonite{} of \NumRegAcidCellsAboveAragonite{} cells from that pH up; lowest passing $\mathrm{Ca_T}$ \NumRegAcidMinCaAragonite{} at pH~\NumRegAcidFirstAragonite{} on the crossed grid; at or above \NumRegHighDose{}) & -- & -- & -- \\
ACC (formation-stage precursor; beside the two, not required) & slurry & suitable against dissolution (\NumRegSlurryPassAcc{} of \NumRegSlurryCellsAcc{} cells; $\mathrm{Ca_T}$ \NumRegWindowLoAcc{} to \NumRegWindowHiAcc{}) & excluded against dissolution\nacreDissolutionMark{} below pH~\NumRegAcidOpenLoAcc{} (lower passivation window, below the sourced range); inconclusive at pH~\NumRegAcidOpenLoAcc{} and \NumRegAcidOpenHiAcc{}; suitable against dissolution from pH~\NumRegAcidFirstAcc{} up (\NumRegAcidPassAcc{} of \NumRegAcidCellsAboveAcc{} cells from that pH up; lowest passing $\mathrm{Ca_T}$ \NumRegAcidMinCaAcc{} at pH~\NumRegAcidFirstAcc{} on the crossed grid; at or above \NumRegHighDose{}) & -- & -- & -- \\
carbonate ($K_1$ control) & control & -- & -- & open & excluded\nacreExcludedMark{} & excluded\nacreExcludedMark{} \\
\midrule
\multicolumn{7}{l}{\emph{slurry-stage phosphates: $K_2$ at the processing stage; $K_1$ is the phase's own ceiling if it were retained, not a gate}} \\
\addlinespace[2pt]
HAp & slurry & suitable against dissolution (\NumRegSlurryPassHap{} of \NumRegSlurryCellsHap{} cells) & suitable against dissolution from pH~\NumRegAcidFirstHap{} (every pH of the grid; \NumRegAcidPassHap{} of \NumRegAcidCellsHap{} cells; lowest passing $\mathrm{Ca_T}$ \NumRegAcidMinCaHap{} and $\mathrm{P_T}$ \NumRegAcidMinPHap{} at pH~\NumRegAcidFirstHap{} on the crossed grid; at its own Ca:P of \NumRatioHap{}, lowest passing $\mathrm{Ca_T}$ \NumRegAcidStoichCaHap{} and $\mathrm{P_T}$ \NumRegAcidStoichPHap{}; both at or above \NumRegHighDose{}) & crystalline: open; amorphous: undetermined & crystalline: excluded\nacreExcludedMark{} at \NumClassMidHigh{}; edge-sensitive at \NumClassMidLow{} (crossover within \NumEdgeTolerance{}); amorphous: excluded (by the bound $\Delta G_{\mathrm{g}\to\mathrm{c}}\ge0$, against the crystalline reference plus water)\nacreExcludedMark{} & crystalline: excluded\nacreExcludedMark{}; amorphous: excluded (by the bound $\Delta G_{\mathrm{g}\to\mathrm{c}}\ge0$, against the crystalline reference plus water)\nacreExcludedMark{} \\
ACP & slurry & suitable against dissolution (\NumRegSlurryPassAcp{} of \NumRegSlurryCellsAcp{} cells) & inconclusive below pH~\NumRegAcidFirstAcp{}; suitable against dissolution from pH~\NumRegAcidFirstAcp{} up (\NumRegAcidPassAcp{} of \NumRegAcidCellsAboveAcp{} cells from that pH up; lowest passing $\mathrm{Ca_T}$ \NumRegAcidMinCaAcp{} and $\mathrm{P_T}$ \NumRegAcidMinPAcp{} at pH~\NumRegAcidFirstAcp{} on the crossed grid; at or above \NumRegHighDose{}) & $\beta$-TCP reference: open; amorphous: undetermined & $\beta$-TCP reference: excluded\nacreExcludedMark{} at \NumClassMidHigh{}; edge-sensitive at \NumClassMidLow{} (crossover within \NumEdgeTolerance{}); amorphous: excluded (by the bound $\Delta G_{\mathrm{g}\to\mathrm{c}}\ge0$, against the crystalline reference plus water)\nacreExcludedMark{} & $\beta$-TCP reference: excluded\nacreExcludedMark{}; amorphous: excluded (by the bound $\Delta G_{\mathrm{g}\to\mathrm{c}}\ge0$, against the crystalline reference plus water)\nacreExcludedMark{} \\
$\beta$-TCP & slurry & suitable against dissolution (\NumRegSlurryPassBtcp{} of \NumRegSlurryCellsBtcp{} cells) & suitable against dissolution from pH~\NumRegAcidFirstBtcp{} (every pH of the grid; \NumRegAcidPassBtcp{} of \NumRegAcidCellsBtcp{} cells; lowest passing $\mathrm{Ca_T}$ \NumRegAcidMinCaBtcp{} and $\mathrm{P_T}$ \NumRegAcidMinPBtcp{} at pH~\NumRegAcidFirstBtcp{} on the crossed grid; at its own Ca:P of \NumRatioTcp{}, lowest passing $\mathrm{Ca_T}$ \NumRegAcidStoichCaBtcp{} and $\mathrm{P_T}$ \NumRegAcidStoichPBtcp{}; both at or above \NumRegHighDose{}) & crystalline: open; amorphous: undetermined & crystalline: excluded\nacreExcludedMark{} at \NumClassMidHigh{}; edge-sensitive at \NumClassMidLow{} (crossover within \NumEdgeTolerance{}); amorphous: excluded (by the bound $\Delta G_{\mathrm{g}\to\mathrm{c}}\ge0$, against the crystalline reference plus water)\nacreExcludedMark{} & crystalline: excluded\nacreExcludedMark{}; amorphous: excluded (by the bound $\Delta G_{\mathrm{g}\to\mathrm{c}}\ge0$, against the crystalline reference plus water)\nacreExcludedMark{} \\
OCP & slurry & inconclusive (control fails) & inconclusive (control fails) & undetermined & undetermined & undetermined \\
brushite & slurry & undetermined (data missing)\textsuperscript{b} & undetermined (data missing) & open & open & excluded above \NumRegCrossBrushiteHi{}\nacreExcludedMark{} \\
$\alpha$-TCP & slurry & undetermined (data missing)\textsuperscript{b} & undetermined (data missing) & crystalline: open; amorphous: undetermined & crystalline: excluded\nacreExcludedMark{}; amorphous: excluded (by the bound $\Delta G_{\mathrm{g}\to\mathrm{c}}\ge0$, against the crystalline reference plus water)\nacreExcludedMark{} & crystalline: excluded\nacreExcludedMark{}; amorphous: excluded (by the bound $\Delta G_{\mathrm{g}\to\mathrm{c}}\ge0$, against the crystalline reference plus water)\nacreExcludedMark{} \\
\ce{Li3PO4} & slurry & undetermined (data missing) & undetermined (data missing) & undetermined & undetermined & undetermined \\
\midrule
\multicolumn{7}{l}{\emph{cell-stage condensed and aluminium phosphates: $K_1$; the route criteria replace $K_2$}} \\
\addlinespace[2pt]
\ce{Ca2P2O7} & cell & route criteria & route criteria & open & open & open \\
\ce{Ca(PO3)2} (crystal) & cell & route criteria & route criteria & open & open & open \\
\ce{Ca(PO3)2} (glass) & cell & route criteria & route criteria & open & open & open \\
\ce{AlPO4} & cell & route criteria & route criteria & undetermined & undetermined & undetermined \\
\ce{Al(PO3)3} & cell & route criteria & route criteria & undetermined & undetermined & undetermined \\
Al-polyphosphate gel & cell & route criteria & route criteria & undetermined & undetermined & undetermined \\
\midrule
\multicolumn{7}{l}{\emph{baselines: $K_1$ and $K_2$; $K_4$ is undetermined for both, for want of a film thickness}} \\
\addlinespace[2pt]
(\ce{LiPO3})$_n$ (baseline) & cell & undetermined (data missing) & -- & open & open & open \\
baked aluminophosphate (baseline) & cell & undetermined (data missing) & -- & undetermined & undetermined & undetermined \\
\bottomrule
\end{tabular}
\par\smallskip
\begin{minipage}{\textwidth}\scriptsize
Verdict words are the registered ones: excluded, open, suitable against dissolution, inconclusive and undetermined (data missing). Open means not falsified at the cutoff, with no lower bound on the onset (in the prose: not thermodynamically excluded); it is not a pass. Inconclusive is a $K_2$ word: the cells of a pH neither all pass nor all fail, or the falsifiability control fails. A dash means the test is not applied to the row; route criteria means the row is cell-stage and the route criteria replace $K_2$.\par
\nacreExcludedMark{} \NumExcludedNote. Kinetic stabilisation is not ruled out for phosphates: tier-two class B records exist in \ce{LiPF6} electrolytes, not in fluorine-free ones.\par
The $K_1$ entries of the calcium orthophosphates are on their crystalline framing (for ACP, its crystalline $\beta$-TCP reference). An amorphous cell of a calcium orthophosphate reads excluded by the bound $\Delta G_{\mathrm{g}\to\mathrm{c}}\ge0$ against the crystalline reference plus water where its crystalline row is excluded, and undetermined otherwise.\par
\nacreDissolutionMark{} Excluded against dissolution (thermodynamic, $K_2$: $\Omega_\mathrm{bulk}<1$ at every cell of that pH); not the tier-one oxidation footnote.\par
Suitable against dissolution is a necessary condition, stated for the cells of the registered grid that pass; it does not address the instability of aluminium above its passivation window and is not evidence of protection. Pass counts are given per row, with the windows: \NumRegWindowCaLo{} to \NumRegWindowCaHi{} for $\mathrm{Ca_T}$ and \NumRegWindowPLo{} to \NumRegWindowPHi{} for $\mathrm{P_T}$, crossed independently on one decade grid (\NumRegGridCaCount{} values of each, \NumRegGridPairs{} pairs per cell group); the acid grid spans pH \NumRegAcidPhLow{} to \NumRegAcidPhHigh{}, the sourced acid-adjusted range is \NumAcidSourcedLo{} to \NumAcidSourcedHi{}, and the slurry grid spans pH \NumRegSlurryPhLow{} to \NumRegSlurryPhHigh{}.\par
Doses are the lowest passing $\mathrm{Ca_T}$ and $\mathrm{P_T}$ of the pH; a minimum at or above \NumRegHighDose{} is flagged in the entry. The sourced \ce{H3PO4} doses could not be converted to a concentration (no liquid volume or density is printed; the only derivable figure is \NumDoseMolPerKgLiquid{} of liquid), so the passing doses are comparable to the order of the sourced additive doses, not equal to them. The dose \NumRegHighDose{} is below that figure and \NumRegDoseSameOrder{} is of the same order as it (the units differ and no liquid density is printed, so the comparison is of order only). Where a phase has its own Ca:P (HAp, $\beta$-TCP) the entry also names its lowest passing doses on that ratio line; ACP has no fixed ratio. See \nacreSIref{si:regime}{Regime table detail and later reruns}.\par
Acid-adjusted column: split by pH (alternative b): below the lowest pH of the grid that has a passing cell, a pH reads excluded against dissolution where every cell is falsified and inconclusive where none passes and some cells are inconclusive; suitable against dissolution from the lowest passing pH up. Slurry-stage $K_1$ is the phase's own ceiling and is not a gate for that stage.\par
\textsuperscript{b}~No $K_{\mathrm{sp}}$ key existed for this phase at the freeze; keys added afterwards are not used in the frozen production, and their reruns are reported beside the frozen rows (\nacreSIref{si:regime}{Regime table detail and later reruns}).\par
\end{minipage}
\endgroup
}

\newcommand{\nacreRegimePerPh}{%
\begingroup\footnotesize\setlength{\tabcolsep}{3pt}
\resizebox{\textwidth}{!}{%
\begin{tabular}{lllllllp{0.12\textwidth}p{0.12\textwidth}lllp{0.12\textwidth}}
\toprule
\multicolumn{3}{l}{} & \multicolumn{6}{c}{crossed decade grid} & \multicolumn{4}{c}{ratio-grid run, own Ca:P} \\
\cmidrule(lr){4-9}\cmidrule(lr){10-13}
phase & basis & pH & passing cells & fraction & lowest passing $\mathrm{Ca_T}$ & lowest passing $\mathrm{P_T}$ & Ca:P of the passing cells & lowest dose at or above \NumRegHighDose{} & own Ca:P & lowest passing $\mathrm{Ca_T}$ & lowest passing $\mathrm{P_T}$ & lowest dose at or above \NumRegHighDose{} \\
\midrule
calcite & frozen & \NumRegPhGridFourFive{} & \NumRegPhPassCalciteFourFive{} of \NumRegPhCellsCalcite{} & -- & -- & -- & -- & -- & -- & -- & -- & -- \\
 &  & \NumRegPhGridFiveFive{} & \NumRegPhPassCalciteFiveFive{} of \NumRegPhCellsCalcite{} & -- & -- & -- & -- & -- & -- & -- & -- & -- \\
 &  & \NumRegPhGridSixFive{} & \NumRegPhPassCalciteSixFive{} of \NumRegPhCellsCalcite{} & \NumRegPhFracCalciteSixFive{} & \NumRegPhMinCaCalciteSixFive{} & -- & -- & yes & -- & -- & -- & -- \\
 &  & \NumRegPhGridSevenFive{} & \NumRegPhPassCalciteSevenFive{} of \NumRegPhCellsCalcite{} & \NumRegPhFracCalciteSevenFive{} & \NumRegPhMinCaCalciteSevenFive{} & -- & -- & no & -- & -- & -- & -- \\
 &  & \NumRegPhGridEightFive{} & \NumRegPhPassCalciteEightFive{} of \NumRegPhCellsCalcite{} & \NumRegPhFracCalciteEightFive{} & \NumRegPhMinCaCalciteEightFive{} & -- & -- & no & -- & -- & -- & -- \\
\addlinespace[2pt]
aragonite & frozen & \NumRegPhGridFourFive{} & \NumRegPhPassAragoniteFourFive{} of \NumRegPhCellsAragonite{} & -- & -- & -- & -- & -- & -- & -- & -- & -- \\
 &  & \NumRegPhGridFiveFive{} & \NumRegPhPassAragoniteFiveFive{} of \NumRegPhCellsAragonite{} & -- & -- & -- & -- & -- & -- & -- & -- & -- \\
 &  & \NumRegPhGridSixFive{} & \NumRegPhPassAragoniteSixFive{} of \NumRegPhCellsAragonite{} & \NumRegPhFracAragoniteSixFive{} & \NumRegPhMinCaAragoniteSixFive{} & -- & -- & yes & -- & -- & -- & -- \\
 &  & \NumRegPhGridSevenFive{} & \NumRegPhPassAragoniteSevenFive{} of \NumRegPhCellsAragonite{} & \NumRegPhFracAragoniteSevenFive{} & \NumRegPhMinCaAragoniteSevenFive{} & -- & -- & no & -- & -- & -- & -- \\
 &  & \NumRegPhGridEightFive{} & \NumRegPhPassAragoniteEightFive{} of \NumRegPhCellsAragonite{} & \NumRegPhFracAragoniteEightFive{} & \NumRegPhMinCaAragoniteEightFive{} & -- & -- & no & -- & -- & -- & -- \\
\addlinespace[2pt]
ACC & frozen & \NumRegPhGridFourFive{} & \NumRegPhPassAccFourFive{} of \NumRegPhCellsAcc{} & -- & -- & -- & -- & -- & -- & -- & -- & -- \\
 &  & \NumRegPhGridFiveFive{} & \NumRegPhPassAccFiveFive{} of \NumRegPhCellsAcc{} & -- & -- & -- & -- & -- & -- & -- & -- & -- \\
 &  & \NumRegPhGridSixFive{} & \NumRegPhPassAccSixFive{} of \NumRegPhCellsAcc{} & -- & -- & -- & -- & -- & -- & -- & -- & -- \\
 &  & \NumRegPhGridSevenFive{} & \NumRegPhPassAccSevenFive{} of \NumRegPhCellsAcc{} & -- & -- & -- & -- & -- & -- & -- & -- & -- \\
 &  & \NumRegPhGridEightFive{} & \NumRegPhPassAccEightFive{} of \NumRegPhCellsAcc{} & \NumRegPhFracAccEightFive{} & \NumRegPhMinCaAccEightFive{} & -- & -- & yes & -- & -- & -- & -- \\
\addlinespace[2pt]
HAp & frozen & \NumRegPhGridFourFive{} & \NumRegPhPassHapFourFive{} of \NumRegPhCellsHap{} & \NumRegPhFracHapFourFive{} & \NumRegPhMinCaHapFourFive{} & \NumRegPhMinPHapFourFive{} & \NumRegPhRatioLoHapFourFive{} to \NumRegPhRatioHiHapFourFive{} & yes & \NumRatioHap{} & \NumRegPhStoichCaHapFourFive{} & \NumRegPhStoichPHapFourFive{} & yes \\
 &  & \NumRegPhGridFiveFive{} & \NumRegPhPassHapFiveFive{} of \NumRegPhCellsHap{} & \NumRegPhFracHapFiveFive{} & \NumRegPhMinCaHapFiveFive{} & \NumRegPhMinPHapFiveFive{} & \NumRegPhRatioLoHapFiveFive{} to \NumRegPhRatioHiHapFiveFive{} & yes & \NumRatioHap{} & \NumRegPhStoichCaHapFiveFive{} & \NumRegPhStoichPHapFiveFive{} & yes \\
 &  & \NumRegPhGridSixFive{} & \NumRegPhPassHapSixFive{} of \NumRegPhCellsHap{} & \NumRegPhFracHapSixFive{} & \NumRegPhMinCaHapSixFive{} & \NumRegPhMinPHapSixFive{} & \NumRegPhRatioLoHapSixFive{} to \NumRegPhRatioHiHapSixFive{} & no & \NumRatioHap{} & \NumRegPhStoichCaHapSixFive{} & \NumRegPhStoichPHapSixFive{} & no \\
 &  & \NumRegPhGridSevenFive{} & \NumRegPhPassHapSevenFive{} of \NumRegPhCellsHap{} & \NumRegPhFracHapSevenFive{} & \NumRegPhMinCaHapSevenFive{} & \NumRegPhMinPHapSevenFive{} & \NumRegPhRatioLoHapSevenFive{} to \NumRegPhRatioHiHapSevenFive{} & no & \NumRatioHap{} & \NumRegPhStoichCaHapSevenFive{} & \NumRegPhStoichPHapSevenFive{} & no \\
 &  & \NumRegPhGridEightFive{} & \NumRegPhPassHapEightFive{} of \NumRegPhCellsHap{} & \NumRegPhFracHapEightFive{} & \NumRegPhMinCaHapEightFive{} & \NumRegPhMinPHapEightFive{} & \NumRegPhRatioLoHapEightFive{} to \NumRegPhRatioHiHapEightFive{} & no & \NumRatioHap{} & \NumRegPhStoichCaHapEightFive{} & \NumRegPhStoichPHapEightFive{} & no \\
\addlinespace[2pt]
ACP & frozen & \NumRegPhGridFourFive{} & \NumRegPhPassAcpFourFive{} of \NumRegPhCellsAcp{} & -- & -- & -- & -- & -- & -- & -- & -- & -- \\
 &  & \NumRegPhGridFiveFive{} & \NumRegPhPassAcpFiveFive{} of \NumRegPhCellsAcp{} & \NumRegPhFracAcpFiveFive{} & \NumRegPhMinCaAcpFiveFive{} & \NumRegPhMinPAcpFiveFive{} & \NumRegPhRatioLoAcpFiveFive{} to \NumRegPhRatioHiAcpFiveFive{} & yes & -- & -- & -- & -- \\
 &  & \NumRegPhGridSixFive{} & \NumRegPhPassAcpSixFive{} of \NumRegPhCellsAcp{} & \NumRegPhFracAcpSixFive{} & \NumRegPhMinCaAcpSixFive{} & \NumRegPhMinPAcpSixFive{} & \NumRegPhRatioLoAcpSixFive{} to \NumRegPhRatioHiAcpSixFive{} & yes & -- & -- & -- & -- \\
 &  & \NumRegPhGridSevenFive{} & \NumRegPhPassAcpSevenFive{} of \NumRegPhCellsAcp{} & \NumRegPhFracAcpSevenFive{} & \NumRegPhMinCaAcpSevenFive{} & \NumRegPhMinPAcpSevenFive{} & \NumRegPhRatioLoAcpSevenFive{} to \NumRegPhRatioHiAcpSevenFive{} & yes & -- & -- & -- & -- \\
 &  & \NumRegPhGridEightFive{} & \NumRegPhPassAcpEightFive{} of \NumRegPhCellsAcp{} & \NumRegPhFracAcpEightFive{} & \NumRegPhMinCaAcpEightFive{} & \NumRegPhMinPAcpEightFive{} & \NumRegPhRatioLoAcpEightFive{} to \NumRegPhRatioHiAcpEightFive{} & yes & -- & -- & -- & -- \\
\addlinespace[2pt]
$\beta$-TCP & frozen & \NumRegPhGridFourFive{} & \NumRegPhPassBtcpFourFive{} of \NumRegPhCellsBtcp{} & \NumRegPhFracBtcpFourFive{} & \NumRegPhMinCaBtcpFourFive{} & \NumRegPhMinPBtcpFourFive{} & \NumRegPhRatioLoBtcpFourFive{} to \NumRegPhRatioHiBtcpFourFive{} & yes & \NumRatioTcp{} & \NumRegPhStoichCaBtcpFourFive{} & \NumRegPhStoichPBtcpFourFive{} & yes \\
 &  & \NumRegPhGridFiveFive{} & \NumRegPhPassBtcpFiveFive{} of \NumRegPhCellsBtcp{} & \NumRegPhFracBtcpFiveFive{} & \NumRegPhMinCaBtcpFiveFive{} & \NumRegPhMinPBtcpFiveFive{} & \NumRegPhRatioLoBtcpFiveFive{} to \NumRegPhRatioHiBtcpFiveFive{} & yes & \NumRatioTcp{} & \NumRegPhStoichCaBtcpFiveFive{} & \NumRegPhStoichPBtcpFiveFive{} & yes \\
 &  & \NumRegPhGridSixFive{} & \NumRegPhPassBtcpSixFive{} of \NumRegPhCellsBtcp{} & \NumRegPhFracBtcpSixFive{} & \NumRegPhMinCaBtcpSixFive{} & \NumRegPhMinPBtcpSixFive{} & \NumRegPhRatioLoBtcpSixFive{} to \NumRegPhRatioHiBtcpSixFive{} & no & \NumRatioTcp{} & \NumRegPhStoichCaBtcpSixFive{} & \NumRegPhStoichPBtcpSixFive{} & yes \\
 &  & \NumRegPhGridSevenFive{} & \NumRegPhPassBtcpSevenFive{} of \NumRegPhCellsBtcp{} & \NumRegPhFracBtcpSevenFive{} & \NumRegPhMinCaBtcpSevenFive{} & \NumRegPhMinPBtcpSevenFive{} & \NumRegPhRatioLoBtcpSevenFive{} to \NumRegPhRatioHiBtcpSevenFive{} & no & \NumRatioTcp{} & \NumRegPhStoichCaBtcpSevenFive{} & \NumRegPhStoichPBtcpSevenFive{} & no \\
 &  & \NumRegPhGridEightFive{} & \NumRegPhPassBtcpEightFive{} of \NumRegPhCellsBtcp{} & \NumRegPhFracBtcpEightFive{} & \NumRegPhMinCaBtcpEightFive{} & \NumRegPhMinPBtcpEightFive{} & \NumRegPhRatioLoBtcpEightFive{} to \NumRegPhRatioHiBtcpEightFive{} & no & \NumRatioTcp{} & \NumRegPhStoichCaBtcpEightFive{} & \NumRegPhStoichPBtcpEightFive{} & no \\
\addlinespace[2pt]
OCP, \ce{Ca4H(PO4)3} basis & post-freeze rerun & \NumRegPhGridFourFive{} & \NumRegPhPassOcpFourFive{} of \NumRegPhCellsOcp{} & \NumRegPhFracOcpFourFive{} & \NumRegPhMinCaOcpFourFive{} & \NumRegPhMinPOcpFourFive{} & \NumRegPhRatioLoOcpFourFive{} to \NumRegPhRatioHiOcpFourFive{} & yes & -- & -- & -- & -- \\
 &  & \NumRegPhGridFiveFive{} & \NumRegPhPassOcpFiveFive{} of \NumRegPhCellsOcp{} & \NumRegPhFracOcpFiveFive{} & \NumRegPhMinCaOcpFiveFive{} & \NumRegPhMinPOcpFiveFive{} & \NumRegPhRatioLoOcpFiveFive{} to \NumRegPhRatioHiOcpFiveFive{} & yes & -- & -- & -- & -- \\
 &  & \NumRegPhGridSixFive{} & \NumRegPhPassOcpSixFive{} of \NumRegPhCellsOcp{} & \NumRegPhFracOcpSixFive{} & \NumRegPhMinCaOcpSixFive{} & \NumRegPhMinPOcpSixFive{} & \NumRegPhRatioLoOcpSixFive{} to \NumRegPhRatioHiOcpSixFive{} & no & -- & -- & -- & -- \\
 &  & \NumRegPhGridSevenFive{} & \NumRegPhPassOcpSevenFive{} of \NumRegPhCellsOcp{} & \NumRegPhFracOcpSevenFive{} & \NumRegPhMinCaOcpSevenFive{} & \NumRegPhMinPOcpSevenFive{} & \NumRegPhRatioLoOcpSevenFive{} to \NumRegPhRatioHiOcpSevenFive{} & no & -- & -- & -- & -- \\
 &  & \NumRegPhGridEightFive{} & \NumRegPhPassOcpEightFive{} of \NumRegPhCellsOcp{} & \NumRegPhFracOcpEightFive{} & \NumRegPhMinCaOcpEightFive{} & \NumRegPhMinPOcpEightFive{} & \NumRegPhRatioLoOcpEightFive{} to \NumRegPhRatioHiOcpEightFive{} & no & -- & -- & -- & -- \\
\addlinespace[2pt]
brushite & post-freeze rerun & \NumRegPhGridFourFive{} & \NumRegPhPassBrushiteFourFive{} of \NumRegPhCellsBrushite{} & \NumRegPhFracBrushiteFourFive{} & \NumRegPhMinCaBrushiteFourFive{} & \NumRegPhMinPBrushiteFourFive{} & \NumRegPhRatioLoBrushiteFourFive{} to \NumRegPhRatioHiBrushiteFourFive{} & yes & \NumRatioBrushite{} & \NumRegPhStoichCaBrushiteFourFive{} & \NumRegPhStoichPBrushiteFourFive{} & yes \\
 &  & \NumRegPhGridFiveFive{} & \NumRegPhPassBrushiteFiveFive{} of \NumRegPhCellsBrushite{} & \NumRegPhFracBrushiteFiveFive{} & \NumRegPhMinCaBrushiteFiveFive{} & \NumRegPhMinPBrushiteFiveFive{} & \NumRegPhRatioLoBrushiteFiveFive{} to \NumRegPhRatioHiBrushiteFiveFive{} & no & \NumRatioBrushite{} & \NumRegPhStoichCaBrushiteFiveFive{} & \NumRegPhStoichPBrushiteFiveFive{} & yes \\
 &  & \NumRegPhGridSixFive{} & \NumRegPhPassBrushiteSixFive{} of \NumRegPhCellsBrushite{} & \NumRegPhFracBrushiteSixFive{} & \NumRegPhMinCaBrushiteSixFive{} & \NumRegPhMinPBrushiteSixFive{} & \NumRegPhRatioLoBrushiteSixFive{} to \NumRegPhRatioHiBrushiteSixFive{} & no & \NumRatioBrushite{} & \NumRegPhStoichCaBrushiteSixFive{} & \NumRegPhStoichPBrushiteSixFive{} & yes \\
 &  & \NumRegPhGridSevenFive{} & \NumRegPhPassBrushiteSevenFive{} of \NumRegPhCellsBrushite{} & \NumRegPhFracBrushiteSevenFive{} & \NumRegPhMinCaBrushiteSevenFive{} & \NumRegPhMinPBrushiteSevenFive{} & \NumRegPhRatioLoBrushiteSevenFive{} to \NumRegPhRatioHiBrushiteSevenFive{} & no & \NumRatioBrushite{} & \NumRegPhStoichCaBrushiteSevenFive{} & \NumRegPhStoichPBrushiteSevenFive{} & yes \\
 &  & \NumRegPhGridEightFive{} & \NumRegPhPassBrushiteEightFive{} of \NumRegPhCellsBrushite{} & \NumRegPhFracBrushiteEightFive{} & \NumRegPhMinCaBrushiteEightFive{} & \NumRegPhMinPBrushiteEightFive{} & \NumRegPhRatioLoBrushiteEightFive{} to \NumRegPhRatioHiBrushiteEightFive{} & yes & \NumRatioBrushite{} & \NumRegPhStoichCaBrushiteEightFive{} & \NumRegPhStoichPBrushiteEightFive{} & yes \\
\bottomrule
\end{tabular}}
\par\smallskip
\begin{minipage}{\textwidth}\scriptsize
Acid-adjusted comparator state, wet slurry, dissolution verdict after the control gate. A cell passes when it is suitable against dissolution; doses are in $\mathrm{mol\,dm^{-3}}$. The $\mathrm{Ca_T}$ and $\mathrm{P_T}$ minima are taken separately. Frozen rows are the production run; post-freeze rows use keys added after the freeze (a necessary condition only; frozen and rerun rows agree wherever both exist).\par
Doses are in $\mathrm{mol\,dm^{-3}}$ and are the lowest passing $\mathrm{Ca_T}$ and $\mathrm{P_T}$ of the pH, taken separately, so they need not occur in one cell. A minimum passing dose at or above \NumRegHighDose{} is flagged in the entry; the flag names a threshold and is not a comparison with the sourced doses. The sourced \ce{H3PO4} additions are \NumDoseWtLo{} to \NumDoseWtHi{} of the active material in five studies\cite{tolchard2023insights,chen2024study,xu2022effect,carvalho2016study,kukay2021aqueous} and \NumDoseMmolPerG{} of one \ce{NMC811}; they could not be converted to a concentration because no liquid volume or density is printed, and the only derivable figure is \NumDoseMolPerKgLiquid{} of liquid, so the passing doses are comparable to the order of the sourced additive doses, not equal to them. The dose \NumRegHighDose{} is below that figure and \NumRegDoseSameOrder{} is of the same order as it (the units differ and no liquid density is printed, so the comparison is of order only). The $\mathrm{P_T}$ window of the phosphoric-acid slurry is a gap of the parameter table.\par
The comparator state is the sourced acid-adjusted slurry, pH \NumAcidSourcedLo{} to \NumAcidSourcedHi{}; the grid spans the published passivation window of aluminium (a grade-D statement; the model windows run from \NumAcidModelLo{} to \NumAcidModelHi{}) and is context for it.\par
The left group of columns is the crossed decade grid, where $\mathrm{Ca_T}$ and $\mathrm{P_T}$ are crossed independently (\NumRegGridPairs{} pairs per cell group), so the Ca:P ratios are powers of ten and the apatite ratios \NumRatioTcp{} and \NumRatioHap{} do not occur. The right group is the ratio-grid run, a separate labelled run beside it, with $\mathrm{P_T}$ at $\mathrm{Ca_T}$ divided by the Ca:P of the phase (\NumRatioHap{} for HAp, \NumRatioTcp{} for $\beta$-TCP, \NumRatioBrushite{} for brushite), so the Ca:P of every passing cell is that ratio. ACP and the \ce{Ca4H(PO4)3} basis of OCP were run on both apatite ratios and have no column of their own; the carbonate rows were not run on it. The verdicts of the ratio-grid run are those of the crossed grid.\par
\end{minipage}
\endgroup
}

\newcommand{\nacreRegimeReruns}{%
\begingroup\footnotesize\setlength{\tabcolsep}{3pt}
\begin{tabular}{lllp{0.2\textwidth}p{0.2\textwidth}}
\toprule
phase or key & frozen run & post-freeze rerun & passing cells, alkaline grid (frozen / rerun of all) & passing cells, acid-adjusted grid (frozen / rerun of all) \\
\midrule
calcite & suitable against dissolution & not rerun & \NumRegRerunCalciteAlkalineFrozen{} (not rerun) & \NumRegRerunCalciteAcidFrozen{} (not rerun) \\
aragonite & suitable against dissolution & not rerun & \NumRegRerunAragoniteAlkalineFrozen{} (not rerun) & \NumRegRerunAragoniteAcidFrozen{} (not rerun) \\
ACC & suitable against dissolution & not rerun & \NumRegRerunAccAlkalineFrozen{} (not rerun) & \NumRegRerunAccAcidFrozen{} (not rerun) \\
HAp & suitable against dissolution & suitable against dissolution & \NumRegRerunHapAlkalineFrozen{} / \NumRegRerunHapAlkalineRerun{} of \NumRegRerunHapAlkalineCells{} & \NumRegRerunHapAcidFrozen{} / \NumRegRerunHapAcidRerun{} of \NumRegRerunHapAcidCells{} \\
ACP & suitable against dissolution & suitable against dissolution & \NumRegRerunAcpAlkalineFrozen{} / \NumRegRerunAcpAlkalineRerun{} of \NumRegRerunAcpAlkalineCells{} & \NumRegRerunAcpAcidFrozen{} / \NumRegRerunAcpAcidRerun{} of \NumRegRerunAcpAcidCells{} \\
$\beta$-TCP & suitable against dissolution & suitable against dissolution & \NumRegRerunBtcpAlkalineFrozen{} / \NumRegRerunBtcpAlkalineRerun{} of \NumRegRerunBtcpAlkalineCells{} & \NumRegRerunBtcpAcidFrozen{} / \NumRegRerunBtcpAcidRerun{} of \NumRegRerunBtcpAcidCells{} \\
OCP, original key & inconclusive (control fails) & inconclusive (control fails) & \NumRegRerunOcpAlkalineFrozen{} / \NumRegRerunOcpAlkalineRerun{} of \NumRegRerunOcpAlkalineCells{} & \NumRegRerunOcpAcidFrozen{} / \NumRegRerunOcpAcidRerun{} of \NumRegRerunOcpAcidCells{} \\
OCP, \ce{Ca4H(PO4)3} basis & not in the frozen run & suitable against dissolution & -- / \NumRegRerunOcpCahAlkalineRerun{} of \NumRegRerunOcpCahAlkalineCells{} & -- / \NumRegRerunOcpCahAcidRerun{} of \NumRegRerunOcpCahAcidCells{} \\
brushite & undetermined & suitable against dissolution & \NumRegRerunBrushiteAlkalineFrozen{} / \NumRegRerunBrushiteAlkalineRerun{} of \NumRegRerunBrushiteAlkalineCells{} & \NumRegRerunBrushiteAcidFrozen{} / \NumRegRerunBrushiteAcidRerun{} of \NumRegRerunBrushiteAcidCells{} \\
$\alpha$-TCP & undetermined & undetermined & \NumRegRerunAtcpAlkalineFrozen{} / \NumRegRerunAtcpAlkalineRerun{} of \NumRegRerunAtcpAlkalineCells{} & \NumRegRerunAtcpAcidFrozen{} / \NumRegRerunAtcpAcidRerun{} of \NumRegRerunAtcpAcidCells{} \\
\ce{Li3PO4} & undetermined & undetermined & \NumRegRerunLipoAlkalineFrozen{} / \NumRegRerunLipoAlkalineRerun{} of \NumRegRerunLipoAlkalineCells{} & \NumRegRerunLipoAcidFrozen{} / \NumRegRerunLipoAcidRerun{} of \NumRegRerunLipoAcidCells{} \\
\bottomrule
\end{tabular}
\par\smallskip
\begin{minipage}{\textwidth}\scriptsize
The rerun used the keys added after the freeze (brushite, \ce{Li3PO4} and the \ce{Ca4H(PO4)3} basis of the octacalcium phosphate key) and the rule that the supersaturated control is the grid state of maximum median $\Omega$; it changed \NumCtlCellsMoved{} cells of the original keys. It is a deviation from the frozen production, reported beside the frozen rows and not in place of them; the carbonate rows were not rerun. Open: the control of the original octacalcium phosphate key fails. \ce{Li3PO4} stays undetermined because its key is a concentration product, the chain has no activity coefficients for it, and no \ce{Li+} dose is registered; $\alpha$-TCP stays undetermined (grade D).\par
\end{minipage}
\endgroup
}

\newcommand{\nacreAlSaltTable}{%
\begingroup\footnotesize\setlength{\tabcolsep}{3pt}
\begin{tabular}{>{\raggedright\arraybackslash}p{0.2\textwidth}>{\raggedright\arraybackslash}p{0.2\textwidth}l>{\raggedright\arraybackslash}p{0.38\textwidth}}
\toprule
phase & route & verdict & reason \\
\midrule
\ce{AlPO4} & direct, to the perchlorate & undetermined & no bound on the exchange entropy of the aluminium perchlorate sink \\
\ce{AlPO4} & stepwise, through \ce{Al(PO3)3} & undetermined & the experimental route through \ce{Al(PO3)3} is undetermined, for want of data on \ce{Al(PO3)3} \\
\ce{Al(PO3)3} & direct, to the perchlorate & undetermined & the experimental route through \ce{Al(PO3)3} is undetermined, for want of data on \ce{Al(PO3)3} \\
Al polyphosphate gel & direct, to the perchlorate & undetermined & the experimental route through \ce{Al(PO3)3} is undetermined, for want of data on \ce{Al(PO3)3} \\
baked aluminophosphate (baseline) & direct, to the perchlorate & undetermined & no bound on the exchange entropy of the aluminium perchlorate sink \\
baked aluminophosphate (baseline) & stepwise, through \ce{Al(PO3)3} & undetermined & the experimental route through \ce{Al(PO3)3} is undetermined, for want of data on \ce{Al(PO3)3} \\
\bottomrule
\end{tabular}
\par\smallskip
\begin{minipage}{\textwidth}\scriptsize
Every row of the aluminium-salt ladder is undetermined. The single most useful datum to fetch is the standard Gibbs energy of formation (or entropy) of solid \ce{Al(ClO4)3}; only its enthalpy of formation exists. \NumAlSaltNoMeta{} of the \NumAlSaltRows{} rows also lack data for \ce{Al(PO3)3}.\par
\end{minipage}
\endgroup
}

\newcommand{\nacreRegimeSensNote}{%
The ladder rebuilt at the sensitivity line of \NumClassSensitivity{} is byte-identical to the nominal ladder, so no verdict changes at \NumClassSensitivity{}.
}

\begin{table*}[tp]
\caption{Regime map of the thermodynamic screen, generated from the onset ladder and the outputs of the main run. Dissolution ($K_2$) is read at the processing stage (grid \NumRegGridId) and in the acid-adjusted comparator, the sourced slurries at pH \NumAcidSourcedLo\ to \NumAcidSourcedHi\ on a grid that spans the passivation window of aluminium, oxidation ($K_1$) per cutoff class. Entries use the vocabulary of the screen and carry their stage; continued below.}\label{tab:regime}
\nacreRegimeTablePartA
\end{table*}
\begin{table*}[tp]
\ContinuedFloat
\caption{Regime map of the thermodynamic screen (continued): slurry-stage phosphates.}
\nacreRegimeTablePartB
\end{table*}
\begin{table*}[tp]
\ContinuedFloat
\caption{Regime map of the thermodynamic screen (continued): cell-stage phosphates and baselines, with the notes of the table.}
\nacreRegimeTablePartC
\end{table*}

\subsection{Processing-stage dissolution stability and the acid-adjusted comparator}\label{sec:res-dissolution}
% (2) Why the alkaline slurry favours the calcium minerals, then the one positive statement.
In the alkaline slurry the deprotonated anions \ce{CO3^{2-}} and \ce{PO4^{3-}} become more abundant with pH, so the supersaturation of a calcium mineral grows with pH. % CL-S07
From the sourced acid constants at infinite dilution, \ce{CO3^{2-}} is about \NumSpecCarbonateLow\ of the dissolved carbonate at pH \NumSpecPhLow. % CL-S07
\ce{PO4^{3-}} is about \NumSpecPhosphateLow\ of the phosphate at pH \NumSpecPhLow\ and \NumSpecPhosphateMid\ at pH \NumSpecPhMid. % CL-S07
The one positive statement reads: thermodynamically suitable against dissolution (a necessary condition) for calcite in grid \NumRegGridId\ at $\mathrm{Ca_T}\ge$ \NumRegWindowFullLoCalcite. % CL-S17
The passing region covers pH \NumRegSlurryPhLow\ to \NumRegSlurryPhHigh\ and \NumRegSlurryPassFullCalcite\ cells, with every carbonate, pH, ionic-strength and drying state passing. % CL-S17
\NumRegSlurryPassExtraCalcite\ further cells pass at \NumRegWindowLoCalcite, where the other cells of that dose fail, so the label does not extend to that dose. % CL-S17
The grid carries no phosphate; the phosphate rows assume $\mathrm{P_T}$ on the same decade grid, \NumRegWindowPLo\ to \NumRegWindowPHi, because no slurry phosphate total is sourced. % CL-S17
This is a necessary condition only: it is not evidence of protection and does not address the instability of aluminium itself above its passivation window. % CL-D20
Aragonite passes in \NumRegSlurryPassAragonite\ of \NumRegSlurryCellsAragonite\ cells, at $\mathrm{Ca_T}$ \NumRegWindowLoAragonite\ to \NumRegWindowHiAragonite. % CL-S17
That is dissolution stability only, because its formation is unsupported at slurry pH and room temperature\cite{nahi2022polyamines,hanifa2025effects}. % CL-S17, CL-K15
Amorphous calcium carbonate, the precursor of the formation stage, passes in \NumRegSlurryPassAcc\ cells, at $\mathrm{Ca_T}$ \NumRegWindowLoAcc\ to \NumRegWindowHiAcc; it is reported beside the two polymorphs and is not required. % CL-S17
Alternating aragonite and calcite layers would not be nacre-like, and the screen tests neither layering nor polymorph structure. % CL-L10
Each of the three reports dissolution at the low-dose corner of the grid, so a pass is not an artefact of the test. % CL-S17
Among the calcium phosphates the same test passes hydroxyapatite in \NumRegSlurryPassHap\ of \NumRegSlurryCellsHap\ cells, $\beta$-TCP in \NumRegSlurryPassBtcp\ and amorphous calcium phosphate in \NumRegSlurryPassAcp. % CL-S07
Octacalcium phosphate is inconclusive because its control fails, and brushite, $\alpha$-TCP and \ce{Li3PO4} are undetermined in the main run. % CL-S07
A later rerun used keys added after the main run. % CL-S18
In it, brushite passes in \NumRegRerunBrushiteAlkalineRerun\ of \NumRegRerunBrushiteAlkalineCells\ alkaline cells and octacalcium phosphate on the \ce{Ca4H(PO4)3} basis in \NumRegRerunOcpCahAlkalineRerun. % CL-S18
Both are reported beside the main-run rows (\nacreSIref{si:regime}{Regime table detail and later reruns}). % CL-S18

% (3) The acid-adjusted comparator: the incumbent route, with both minimum doses.
Phosphoric acid in the incumbent route brings the slurry to pH \NumAcidSourcedLo\ to \NumAcidSourcedHi, the sourced acid-adjusted state and the comparator of this section. % CL-S08
The dissolution test runs on a grid from pH \NumRegAcidPhLow\ to \NumRegAcidPhHigh\ that spans the passivation window of aluminium. % CL-S08
That window is a grade-D statement, and our model windows run from \NumAcidModelLo\ to \NumAcidModelHi. % CL-S08
Lowering the pH protonates carbonate and phosphate and drives every calcium mineral toward dissolution. % CL-S08
Calcite and aragonite dissolve at every cell at pH \NumRegAcidExclLoCalcite\ and \NumRegAcidExclHiCalcite. % CL-S08
Calcite passes from pH \NumRegAcidFirstCalcite\ up, from a calcium total of \NumRegAcidMinCaCalcite\ at that pH, in \NumRegAcidPassCalcite\ of the \NumRegAcidCellsAboveCalcite\ cells at that pH and above, inside the sourced range. % CL-S08
The grid has no point between these pH values. % CL-S08
Carbonate is therefore excluded against dissolution at the two lowest pH values of the grid, in the lower part of the passivation window. % CL-S08
It is not excluded against dissolution at the grid points inside the sourced range (pH \NumRegPhGridSixFive, \NumRegPhGridSevenFive\ and \NumRegPhGridEightFive). % CL-S08
Hydroxyapatite and $\beta$-TCP pass at every pH of the acid-adjusted grid, but at pH \NumRegAcidFirstHap\ only at high dose. % CL-S08
There, hydroxyapatite passes from $\mathrm{Ca_T}$ \NumRegAcidMinCaHap\ and $\mathrm{P_T}$ \NumRegAcidMinPHap\ on the crossed decade grid, and from \NumRegAcidStoichCaHap\ and \NumRegAcidStoichPHap\ at its own Ca:P of \NumRatioHap. % CL-S08
$\beta$-TCP passes from \NumRegAcidMinCaBtcp\ and \NumRegAcidMinPBtcp\ on the crossed grid, and from \NumRegAcidStoichCaBtcp\ and \NumRegAcidStoichPBtcp\ at \NumRatioTcp. % CL-S08
At their own stoichiometry both pass at low dose at higher pH. % CL-S08
Hydroxyapatite passes from $\mathrm{Ca_T}$ \NumRegOwnHapAcidCa\ ($\mathrm{P_T}$ \NumRegOwnHapAcidP) at pH \NumRegOwnHapAcidPh, and $\beta$-TCP from \NumRegOwnBtcpAcidCa\ (\NumRegOwnBtcpAcidP) at pH \NumRegOwnBtcpAcidPh. % CL-S08
In the later rerun, brushite stays a high-dose phase at pH \NumRegPhGridFourFive, needing $\mathrm{Ca_T}$ \NumRegPhStoichCaBrushiteFourFive\,$\mathrm{mol\,dm^{-3}}$ on its own ratio line. % CL-S08, CL-S18
Octacalcium phosphate, read on the \ce{Ca4H(PO4)3} basis, passes in \NumRegRerunOcpCahAcidRerun\ of \NumRegRerunOcpCahAcidCells\ cells in the later rerun. % CL-S18
$\alpha$-TCP (grade D) and \ce{Li3PO4} (a concentration-product key) remain undetermined. % CL-S08
These passing doses are comparable to the order of the sourced additive doses, never equal to them. % CL-S08
The one derivable sourced figure is about \NumDoseMolPerKgLiquid\ of liquid, in other units. % CL-S08
The table flags passing doses at or above \NumRegHighDose, of which \NumRegDoseSameOrder\ is of the same order as that figure. % CL-S08
The sourced acid-adjusted state is the incumbent comparator, not a regime of the barrier, and the grid that spans the passivation window is context for it. % CL-S08, CL-M13

\subsection{The carbonate control and the $K_1$ exclusion map}\label{sec:res-carbonate}
% (4) The physics of the oxidation onset of an isolated-oxoanion mineral, then the control.
% carbonate control (frozen)
At the cutoff a calcium mineral oxidises once releasing its calcium to the electrolyte salt and its lattice oxygen as oxygen gas, with the electrons plating lithium, lowers the free energy. % CL-S03
Taking the solid perchlorate as the sink, with the lithium salt drawn from its solid, gives an upper bound on that onset for the saturated sink\nacreSinkMark. % CL-S03
Below saturation the bound rises by the one-sided term $m_\mathrm{sink}$ of Section~\ref{sec:design}, whose value is not evaluated. % CL-S03
For calcium carbonate that bound is \NumCarbonateOnset. % CL-S03
With its margin, carbonate is not thermodynamically excluded at the nominal low class of \NumClassLo. % CL-S03
It is thermodynamically excluded\nacreExcludedMark\nacreSinkMark\ above \NumCarbonateCrossover, which is every cutoff of the two higher classes (Figs.~\ref{fig:ladder} and~\ref{fig:heatmap}, where the carbonate $K_1$ cells are those of calcite). % CL-S03
Lithium carbonate, oxidised in glyme electrolytes, evolves gas from \NumTierTwoLow\ to \NumTierTwoHigh, close to the equilibrium value of its own reaction\cite{mahne2018electrochemical,cao2022oxidative}. % CL-S15
Read across to calcium carbonate, this is a bulk-phase record of another phase, so the claim is worded as low voltage up to about \NumTierTwoLow\ (a contrary reading is in \nacreSIref{si:templating}{Templating detail}). % CL-S15
At the \NumClassSensitivity\ sensitivity line the carbonate crossover lies within the margin of the line, so carbonate is close to thermodynamic exclusion there. % CL-S15
Rebuilt at that line, the ladder changes no verdict; the files are identical, because class verdicts use the class bounds. % CL-S15
% (5) Bone mineral beside carbonate, and the rest of the ladder.
Hydroxyapatite and the tricalcium phosphates strip calcium from a lattice of isolated anions in the same way, and their upper bounds lie at \NumOnsetBoneLow\ to \NumOnsetBoneHigh, beside that of carbonate. % CL-S01
Hydroxyapatite, $\beta$-TCP and amorphous calcium phosphate on its crystalline $\beta$-TCP reference are thermodynamically excluded\nacreExcludedMark\nacreSinkMark\ at \NumClassMidHigh\ and above, and $\alpha$-TCP from \NumClassMidLow. % CL-S14
The crossovers of hydroxyapatite (\NumHapCrossover) and $\beta$-TCP (\NumBtcpCrossover) lie within \NumEdgeTolerance\ of the \NumClassMidLow\ edge, so their exclusion at that edge depends on the sink term. % CL-S14
On the amorphous framing these phases are thermodynamically excluded\nacreExcludedMark\ with them, at the same classes. % CL-S14
The reason is the bound $\Delta G_{\mathrm{g}\to\mathrm{c}}\ge0$ against the crystalline reference of the same composition plus liquid water at unit activity; amorphous calcium phosphate is metastable against crystalline calcium phosphates in water. % CL-S14
Hydroxyapatite is therefore excluded under both framings at \NumClassMidHigh\ and above, and at the \NumClassMidLow\ class within the edge sensitivity stated above. % CL-S14
At the lowest class (\NumClassLo) the crystalline rows are open and the amorphous framing of every phosphate stays undetermined. % CL-S14
The bound does not cover the amorphous framing of every other phosphate, nor amorphous calcium carbonate. % CL-S14
Brushite is excluded only in part: it is not thermodynamically excluded at the two lower classes. % CL-S14
It is thermodynamically excluded\nacreExcludedMark\nacreSinkMark\ above \NumRegCrossBrushiteHi\ inside the highest class, so at \NumClassHiHigh, by oxidation of its hydrate water at the water activity fixed for the screen. % CL-S14
Brushite is a precursor, not a cell-stage phase, and its crossover lies within \NumEdgeTolerance\ of the \NumClassHiLow\ edge. % CL-S14
Table~\ref{tab:breakeven} gives the headroom of each exclusion\nacreSinkMark\ below its cutoff and its break-even ratio $r^*$. % CL-S03, CL-S14
Read plainly, an exclusion stands while the mean-activity ratio of saturated to actual \ce{LiClO4} stays below $r^*$. % CL-S03, CL-S14
Octacalcium phosphate, \ce{Li3PO4}, \ce{Al(PO3)3} and the aluminium polyphosphate gel have no oxidation bound and are omitted from the map. % CL-S12
The silicates \ce{CaSiO3}, \ce{Ca2SiO4} and \ce{Li2SiO3} are comparators of the oxidation ladder, not candidate barriers: calcium and lithium are held by silicate instead of phosphate anions. % CL-S20
\ce{CaSiO3} and \ce{Li2SiO3} are undetermined, because no uncertainty is printed for \ce{SiO2} (\nacreSIref{si:thermo}{Thermochemistry benchmark and the DFT spread}). % CL-S20
\ce{Ca2SiO4} is thermodynamically excluded\nacreExcludedMark\nacreSinkMark\ already at \NumClassLo. % CL-S20
Its limiting oxidation hands half its calcium to the perchlorate sink and leaves \ce{CaSiO3}, with an upper bound of \NumBenCaSilOrthoROneSEub\ and a margin of \NumBenCaSilOrthoROneSMargin. % CL-S20
Their sum lies below that class (Figs.~\ref{fig:ladder} and~\ref{fig:heatmap}, Table~\ref{tab:breakeven}). % CL-S20
The lower end of every upper bound is conditional on the transfer premise, which is pending, so the bars of Fig.~\ref{fig:ladder} are open below. % CL-M17
These exclusions are thermodynamic. % CL-L09
Kinetic stabilisation is not ruled out for phosphates, because tier-two records of coated cathode particles exist in \ce{LiPF6} electrolytes\cite{xiao2015unravelling,hofmann2020surface}, not in fluorine-free ones. % CL-L09
The upper bounds of this ladder were seen before the regime framing of the paper was adopted. % CL-M08

\begin{figure*}[t]
  \centering
  \includegraphics[width=\textwidth]{\nacreRepoRoot/figures/out/fig17.pdf}
  \caption{(a) $K_2$ stability map by slurry pH, rows from the bottom by $K_1$ upper bound $E_\mathrm{ub}$ (calcite, $\alpha$-TCP, ACP, $\beta$-TCP, HAp, brushite). Circle area: passing fraction, filled below a lowest passing dose of \NumRegHighDose, open above; cross: excluded ($K_2$: every cell dissolves; $K_1$: thermodynamically excluded$^\dagger$); plus: open; dot: undetermined. (b) $K_1$ upper bounds only, bars open below; crossed hatching, dashed edges: incomplete budget; grey: cutoff classes; \NumRegNoKOneBound\ phases without a bound omitted.}
  \label{fig:ladder}
\end{figure*}

\begin{figure*}[t]
  \centering
  \includegraphics[width=\textwidth]{\nacreRepoRoot/figures/out/fig15.pdf}
  \caption{Regime of each candidate (rows: slurry phases, cell phases, silicates, carbonates; columns: $K_2$ in the processing and acid-adjusted slurries by pH, then $K_1$ at the cutoff classes, nominal \NumClassLo, crystalline left, amorphous right). Dots: suitable, vertical lines at a lowest passing dose of at least \NumRegHighDose; plain: undetermined; single hatching: open ($K_1$) or inconclusive ($K_2$); cross: excluded ($K_2$: every cell dissolves; $K_1$: thermodynamically excluded$^\dagger$). Brushite at \NumClassHiLow--\NumClassHiHigh\ is thermodynamically excluded only above \NumRegCrossBrushiteHi.}
  \label{fig:heatmap}
\end{figure*}

\subsection{Phosphates and the condensed and aluminium phosphates}\label{sec:res-phosphate}\label{sec:res-condensed}
% (6) Why condensation raises the onset; the amorphous framing; the aluminium phases.
For hydroxyapatite the limiting oxidation releases lattice oxygen and condenses the remaining phosphate to pyrophosphate. % CL-S05
That reaction lies below the oxidation of its hydroxide, so linking phosphate into chains is itself the reaction that sets the ceiling. % CL-S05
A phase whose phosphate is already condensed can release lattice oxygen only by condensing further toward the phosphorus oxide, and its upper bound rises. % CL-S06
The bound is \NumPyroOnset\ for \ce{Ca2P2O7} and \NumMetaGlassOnset\ for \ce{Ca(PO3)2} glass, against \NumMetaCrystalOnset\ for its crystal. % CL-S06
With no lower bound these phases are not thermodynamically excluded at any cutoff class, which is not a pass. % CL-S06
The in-situ routes give glasses, and a glass lies above its crystal in free energy, so the glass row governs and the crystal is an optimistic reference. % CL-S06
The amorphous row of \ce{Ca2P2O7} is undetermined for want of glass data. % CL-S06
The aluminium phosphates, \ce{AlPO4}, \ce{Al(PO3)3}, the aluminium polyphosphate gel and the baked aluminophosphate, are undetermined (\nacreSIref{si:regime}{Regime table detail and later reruns}). % CL-S12
The exchange entropy of the aluminium perchlorate sink has no bound, and \ce{Al(PO3)3} has no thermochemistry. % CL-S12
The single most useful datum for that ladder is the standard Gibbs energy of formation, or the entropy, of solid \ce{Al(ClO4)3}, whose enthalpy alone is tabulated. % CL-S12
% (7) The database benchmark as a method finding.
Database formation energies differ from tabulated enthalpies by \NumDftCaPhosLow\ to \NumDftCaPhosHigh\ per formula unit over the calcium phosphates of the onset ladder. % CL-X03
The cause is a method finding: the \NumMpCorrectionScheme\ scheme applies one oxide correction, \NumMpOxideCorr\ per oxygen atom fitted on ionic oxides, to every oxygen. % CL-X03
That includes lattice water, where the database error of brushite differs from that of monetite by \NumBrushiteWaterExcess\ per lattice water, and the covalent oxygen of condensed phosphates. % CL-X03
On one reaction the density functional onset differs from the experimental route by \NumDftOnsetSpreadLow\ to \NumDftOnsetSpreadHigh\ (density functional minus experiment, \NumDftOnsetPairs\ reaction-scheme pairs, median \NumDftOnsetMedian). % CL-X03
Density functional values therefore appear only as labelled comparators and decide no row (\nacreSIref{si:thermo}{Thermochemistry benchmark and the DFT spread}). % CL-X03
% (8) The route layer, screened as thermodynamics.
A condensed film reaches the cell only by a route that forms it before its chains hydrolyse in the slurry, and screened as thermodynamics every route is undetermined. % CL-S09
The hydrolysis descriptor is outside its data at every process state: no temperature tolerance, transfer factor or wet residence time is sourced. % CL-S09
Gelation at the oxide lacks a solubility product of the aluminium polyphosphate, and the coacervate route lacks a phase map. % CL-S09
The race of gelation or drying against hydrolysis stays a labelled hypothesis. % CL-S09, CL-D13

\subsection{Nucleation and templating as model statements}\label{sec:res-nucleation}
% (9) The F3 sentence form, numbers from the file.
For a solid phosphate nucleus on a hydrophilic template, a contact angle below 90 degrees requires $\gsn$ below $\gsl$, which lies between \NumNucGslQFive\ and \NumNucGslQNinetyFive\ (median \NumNucGslMedian; admitted draws span \NumNucGslMin\ to \NumNucGslMax). % CL-T03
The ignorance prior puts \NumNucMassSourced\ of its mass there under the sourced calcite $\gamma_\mathrm{nv}$, and at most \NumNucMassDesignMax\ over the design points with enough admitted draws, whatever $\gamma_\mathrm{nv}$ (its support widens with $\gamma_\mathrm{nv}$). % CL-T03
A measured or computed $\gsn$ to within about \NumNucDecideWithin\ would decide it. % CL-P01
The nucleation half of the model failed its retrodiction (\NumRetroCorrect\ of \NumRetroUsableCases\ against \NumRetroPassCount), so this is a model statement, not evidence that templated nucleation is impossible. % CL-T03
% (10) The carbonate chain, the deletion test and the sensitivity ranking.
The carbonate chain is ruled out under the model's assumptions at \NumNucLinkFalsified\ of \NumNucLinkDesigns\ design points; the other \NumNucLinkInconclusive\ are inconclusive. % CL-T03
The deletion test asks whether a pass degrades when the template is replaced by the bare oxide. % CL-S19
It finds no degradation at any of the \NumDeletionFramings\ point-framings of the carbonate chain, because neither the template nor the bare oxide passes, so no step of it can be credited to the template. % CL-S19
The Sobol indices of the nucleation link reach a half-width of \NumSobolNucHalfWidth\ against the target of \NumSobolNucTarget, so the Shapley effects carry the ranking. % CL-S19
In them, under the stated priors, the template--nucleus energy governs the shape factor, with a Shapley share of \NumShapleyGsnShape. % CL-S19
The nucleation-window and thin-film conditions of every phosphate framing are undetermined, because the nucleus--liquid and nucleus--vapour energies, the molar volumes or the hydration numbers are gaps. % CL-S19

\subsection{The required-precision map}\label{sec:res-precision}
% (11) What would decide the open verdicts.
Each open verdict has a flip value, and the precision a measurement needs to cross it is the useful output where the screen leaves a verdict open. % CL-P04
The map has \NumPrecRows\ rows (\nacreSIref{si:sensitivity}{Sensitivity detail}). % CL-P04
\NumPrecClosed\ are closed, decided at the prior or by a computed bound; \NumPrecMonteCarlo\ are set by Monte Carlo contraction of the prior; and \NumPrecGap\ are gaps with no prior to contract. % CL-P04
Two measurements decide the most. % CL-P04
A pass of the nucleation condition would need $\gsn$ of hydroxyapatite within about \NumNucDecideWithin\ of zero, beyond current methods. % CL-P04
Confirming its failure needs a precision of only about \NumPrecGammaSnConfirm, which is within reach. % CL-P04
The relative transfer free energy $\Delta G_\mathrm{tr}^\mathrm{rel}$ is that of \ce{Ca^2+} less that of two \ce{Li+}, from water to propylene carbonate. % CL-P04
Measured to about \NumPrecTransferCaKj, a charge-balanced headroom and not a single-ion precision, it decides the low class for the calcium orthophosphates. % CL-P04
The \ce{Ca(PO3)2} glass could then pass up to the highest class (headroom \NumPrecMetaGlassHiHeadroom), and \ce{Ca2P2O7} passes no class above the low class whatever the transfer data. % CL-P04
For the adhesion descriptor, the two contact energies of the Dupr\'e comparator would have to be measured together to about \NumPrecDupreMj. % CL-P04
For the aluminium-salt ladder the single most useful datum is the Gibbs energy or entropy of solid \ce{Al(ClO4)3}. % CL-S12

\subsection{The baselines}\label{sec:res-baselines}
% (12) The two baselines (comparator films, named by chemistry), through K1 and K2 only.
The two baselines, a lithium polyphosphate film and a baked aluminophosphate, enter only through the oxidation, dissolution and thin-film rows, and no ranking against them follows. % CL-S11
The thin-film row is undetermined for both, for want of a film thickness. % CL-S11
The lithium polyphosphate film is not thermodynamically excluded at any cutoff class on its glass row, with an upper bound of \NumLipoGlassOnset. % CL-S11
Its dissolution row is undetermined, since no solubility for its composition is sourced and its water solubility is an inference from chemistry. % CL-S11
The baked aluminophosphate is undetermined in the oxidation and dissolution rows, for want of the aluminium-salt datum above and of a solubility for its composition. % CL-S11
